\chapter{核心模块详细设计与实现}
\label{chap:modules}

本章结合 \texttt{src/} 实现，对核心模块作详细说明。

\section{ECS 核心模块}

\texttt{EntityManager} 使用自增 ID 与空闲列表复用实体。\texttt{ComponentStore} 按类型维护 \texttt{Map<EntityId, Component>}。\texttt{EcsWorld.update} 依次调用已注册系统，按 \texttt{input/logic/physics/render} 分组排序，与 Gregory 所述引擎主循环驱动子系统更新的模式相对应\cite{gregory2018game}。

\texttt{ViewComponent}、\texttt{BodyUIComponent} 绑定 Laya 节点；\texttt{ViewSyncSystem} 同步坐标；\texttt{BloodBarSyncSystem} 用 hp/maxHp 更新 ProgressBar。

\section{属性、移动与筛选器}

\texttt{AttributeSystem} 支持 modifier 按来源增删，供 Buff 与升级奖励使用\cite{briand2001maintainability}。\texttt{MovementSystem}、\texttt{ControlSystem} 处理位移与输入。\texttt{FilterRegistry} 提供 \texttt{Players}、\texttt{Monsters} 等命名筛选，减少重复查询。

\section{技能与效果执行}

\texttt{PlayerAutoCastSystem} 对三装备栏维护独立冷却与 \texttt{pendingCasts}。\texttt{EffectExecutor} 将配表 effect 映射为子弹、穿透/分裂/连锁修饰或直伤。修饰器须位于后续 \texttt{bullet} effect 之前，形成链式组合语义。

\begin{table}[htbp]
  \centering
  \caption{效果类型与执行器映射}
  \label{tab:effect-map}
  \begin{tabular}{ll}
    \toprule
    \texttt{effect} 字段 & 行为 \\
    \midrule
    \texttt{bullet} & \texttt{BulletSystem.spawnBulletWithOptions} \\
    \texttt{modifier\_split} & 设置分裂数量 \\
    \texttt{modifier\_chain} & 设置连锁与半径 \\
    \texttt{modifier\_pierce} & 增加穿透 \\
    \texttt{direct\_damage} & 公式直伤 \\
    \bottomrule
  \end{tabular}
\end{table}

\section{子弹系统与碰撞}

\texttt{BulletSystem} 从 \texttt{BulletPool} 取节点，更新运动，用 \texttt{hitRadiusSq} 做圆碰撞，按阵营过滤后扣血并处理穿透。密集命中场景可参考人群仿真中的局部相互作用与压力模型思想\cite{helbing2000panic,reynolds1987flocks}，本项目采用简化半径检测以降低开销。

\section{怪物系统}

\texttt{MonsterWaveSpawnSystem} 在玩家周围环带刷怪；\texttt{MonsterChaseSystem} 追逐并叠加分离力与随机摆动，Worker 可参与数值重算\cite{cantrell2018unity}。\texttt{MonsterContactDamageSystem} 处理接触伤害；\texttt{MonsterRecycleSystem} 与 \texttt{MonsterPool} 回收离屏单位。

\section{进度、元游戏与跑图}

\texttt{ExperienceSystem}、\texttt{UpgradeRewardSystem} 与奖励面板构成成长闭环。\texttt{MetaFlowController} 协调菜单与 \texttt{onEnterCombat}。\texttt{RunMapGenerator} 生成三幕 DAG，校验 Boss 可达，失败时使用保底图。

\section{技能装配 UI}

\texttt{SkillSelectPanelController} 以 Tab 切换面板并设置 \texttt{GameSession.paused}。\texttt{SkillDragService} 长按拖拽，\texttt{SkillLoadoutSyncSystem} 写回战斗实体。\texttt{UIStackManager} 管理全屏路由，生命周期由 \texttt{UiControllerBase} 约定。

\section{配置加载}

\texttt{ConfigBootstrap} 双通道加载 JSON 并 \texttt{initData} 填充 \texttt{Data}，是数据驱动枢纽\cite{gregory2018game,munaiah2017github}。

\section{典型场景剖析}

\textbf{首次进入战斗：} \texttt{Main} 加载配表 $\rightarrow$ 元菜单选关 $\rightarrow$ \texttt{demo.init()} 创建玩家与系统 $\rightarrow$ 生存计时启动。

\textbf{Tab 装配技能：} 暂停战斗 $\rightarrow$ 拖拽更新 \texttt{SkillLoadoutState} $\rightarrow$ 同步至自动施法 $\rightarrow$ 恢复战斗。

\textbf{高密度战斗：} 实体数达标时 \texttt{CombatDataBridge} 启用 Worker，主线程仍处理子弹碰撞与渲染\cite{cantrell2018unity,bott2014unreal}。

\section{本章小结}

本章详述了 ECS、技能、子弹、怪物、元游戏、UI 与配置等模块实现，下一章讨论性能优化。
